\section{Licences, release and AI use}\label{app:licences}

We list the licence of every upstream piece, say how our files are licensed, and state the release status as last checked. We make no claim about prize eligibility. The licence survey is from 2026-10-13 and was not re-fetched for this version.

\paragraph{Upstream licences.} Table~\ref{tab:licences} gives what we found on 2026-10-13 on the model cards, repositories and Kaggle pages. Readings of ours are marked. This is not legal advice.

\begingroup\footnotesize
\setlength{\tabcolsep}{3pt}
\renewcommand{\arraystretch}{1.1}
\begin{longtable}{@{}>{\raggedright\arraybackslash}p{0.22\linewidth}>{\raggedright\arraybackslash}p{0.45\linewidth}>{\raggedright\arraybackslash}p{0.25\linewidth}@{}}
\caption{Licence of every upstream piece, as found on 2026-10-13. ``Our reading'' marks an inference, not a finding.}\label{tab:licences}\\
\toprule
Piece & Licence found & Note \\
\midrule
\endfirsthead
\toprule
Piece & Licence found & Note \\
\midrule
\endhead
\bottomrule
\endfoot
Qwen3.8-Flash-Next \citep{qwen38flashnext}
& Qwen Community License 1.0 (text of the LICENSE file in the model repository): an MIT-style grant (use, modify, distribute, sublicense, fine-tune, derivatives) with conditions. The copyright and permission notice must travel with copies. The model name must be shown on products above 100M monthly users or US\$20M monthly revenue. A separate Qwen licence is needed to run a ``Model as a Service'' or ``AI Work Assistant'' business (internal use excepted). % N-E6s6.2 audit_cites-C25
& Our reading: a custom licence with field-of-use conditions, not an OSI-approved text. \\
\addlinespace
Intel W4A16 AutoRound \citep{intel2026w4a16}
& Model card ``other / qwen-community-1.0'', ``follow the license of the original model''; Kaggle copy ``Other (specified in description)''.
& Derivative of the Qwen weights. \\
\addlinespace
albucino drafter \citep{albucino2026drafter}
& \texttt{qwen-community-1.0}; Kaggle copy ``Other (specified in description)''.
& \\
\addlinespace
ursuciprian D3 weights \citep{ursuciprian2026d3}
& \texttt{qwen-community-1.0}, ``the license conditions apply to this derivative''; training tools Apache-2.0.
& Our public dataset (version 2) holds the refit tensors unchanged and so inherits this (our reading); version 2 carries the Qwen licence text and a NOTICE with this attribution. \\ % N-E6s6.2 V3-W42
\addlinespace
lee-chang-93 REAP-448 checkpoint; REAP method \citep{lasby2025reap}
& Checkpoint: NVIDIA Open Model License and Qwen Community License 1.0. Method repository: Apache-2.0.
& We use only which expert ids were kept, never their weights. \\
\addlinespace
Pennyroyal SGLang, SGLang, Ratsimbazafy fork
& Apache-2.0.
& \\
\addlinespace
Olympie fork (low-M GEMM patch)
& No licence file.
& An upstream problem in Franzen's wheelhouse, not ours. \\
\addlinespace
lovedheart sglang fork (QSA fix)
& Apache-2.0.
& \\
\addlinespace
Tufa Duck harness \citep{bessis2026duck}
& No LICENSE file at the repository root, in \texttt{ARC3-Inference} or in \texttt{tufa-arc-agi-framework}. The \texttt{pyproject.toml} classifier reads ``License :: OSI Approved :: MIT License''.
& A classifier is a claim, not a grant. \\
\addlinespace
Franzen's notebook and patch \citep{franzen2026milestone2}
& His repository is Apache-2.0. The notebook is a public Kaggle notebook; our fork's first cell states ``Apache 2.0 by default'' and keeps the original's terms.
& \\
\addlinespace
Franzen's two datasets
& Kaggle licence field ``Unknown'' for both; both public.
& The forum asks whether rule 5.a.3 covers ``unknown''; no answer. \\
\addlinespace
shiiin9 D-prime file \citep{shiiin92026dprime}
& The source notebook is public. Its licence is not recorded in our repository and we did not read it.
& Kaggle's default would be Apache-2.0 (our reading). \\
\addlinespace
ARC-AGI repository (arcprize)
& MIT.
& The competition wheels are separate artifacts. \\
\addlinespace
Our D3 dataset and kernel
& Dataset: public, version 2, with the Qwen Community License 1.0 text and a NOTICE (ursuciprian attribution); our 2026-10-13 survey found no licence set. Kernel: private on 2026-11-02; not re-checked since.
& See below. \\
\end{longtable}
\endgroup

\paragraph{Our files.} We will release the files we wrote under CC-BY 4.0, the licence that Kaggle's competition rules apply to a winning submission and its source code; third-party files keep their own licences (Table~\ref{tab:release}). The 12 overlays behind the 40.98 configuration hold 1,316 lines, and ten of them (988 lines) are ours alone. % N-E6s6.4 H9
The ARC Prize page separately asks that code authored by the submitter be released under a permissive public-domain licence such as CC0 or MIT-0 (below); the two texts differ, so our code is offered under either CC BY 4.0 or MIT-0 at the reader's choice, and our paper, figures and data under CC BY 4.0. The submitted kernel, version 1, cannot be edited without a new version, and its first cell says Apache 2.0 by default. Copyright in AI-generated code is unsettled in some jurisdictions, and we have not checked any assistant's terms on output ownership.

\begin{table}[!htb]
\caption{How our files and the files inside them are licensed. Lines are file lines including docstrings, except the last row (added lines).}\label{tab:release}
\centering\footnotesize
\setlength{\tabcolsep}{3pt}
\renewcommand{\arraystretch}{1.1}
\begin{tabular}{@{}>{\raggedright\arraybackslash}p{0.40\linewidth}r>{\raggedright\arraybackslash}p{0.44\linewidth}@{}}
\toprule
Files & Lines & Licence \\
\midrule
wastefix, replycache, agentfix, memcap, toolfast, promptobs, watchdog, failclosed, asyncmrope\_serve, d3draft\_serve & 988 & CC-BY 4.0 for our lines; the anchor strings they quote keep their owners' licences. \\ % N-E6s6.4
p448\_serve & 214 & CC-BY 4.0 for the code; the table of kept expert ids is derived from third-party models, so credit it. \\ % N-E6s6.4
qsaring\_serve & 114 & Stays Apache-2.0 with its notice (derived from Apache-2.0 code). \\ % N-E6s6.4
Build script and tests & not counted & CC-BY 4.0. \\
shiiin9's D-prime file & 252 & Not ours; we cannot relicense it. \\ % N-E6s6.4
Franzen's patch, Duck harness, SGLang, weights & 8,405 & Third-party; cannot be dedicated by us. \\ % N-E6s6.4
\bottomrule
\end{tabular}
\end{table}

\paragraph{Release status (as of 2026-11-03).} The repository is public at github.com/vineethsaivs/arc3-decisions-or-tokens (tag v1, commit c77243c) with a LICENSE file and a THIRD\_PARTY list. It holds the paper sources and figures, the analysis scripts and data, the overlays and build script (without the files whose licences we could not confirm), the gate records and the research notes; the figure of the noise ruler and the discovery figure rebuild from it with one command. It leaves out the Tufa Duck bundle, shiiin9's priority file and the expert-id table, so most overlay tests do not run from it alone. The D3 dataset is public as version 2, which carries Qwen's licence text and a NOTICE with ursuciprian's attribution (reported 2026-11-03). The submitted kernel was private on 2026-11-02 and we have not re-checked it: it embeds those files as one escaped string, so a reader cannot audit the code from the notebook. % V3-W40 V3-W42 V3-W42
\paragraph{Release record.} Repository github.com/vineethsaivs/arc3-decisions-or-tokens, tag v1, commit c77243c, public from 2026-11-02; anonymous requests for the repository and its LICENSE returned 200 that day. The dataset (version 2) became public afterwards; the kernel was not yet public on 2026-11-02 and is to be re-checked from a logged-out browser.
\ownernote{Before posting: confirm the kernel is public (the D3 dataset already is), push the repository at a tagged commit with a LICENSE file, the overlays and the build script, then rewrite the release record above and replace the status paragraph with ``Released at <repo>@<hash> on <date>''.}
The dataset carries Qwen's licence text and ursuciprian's attribution (version 2). Franzen's four inputs are public today. If one goes private, every submission errors, so a fallback build mounts our copies from a private kernel that would also need to be public to be cited.

\paragraph{Rules.} The ARC Prize 2026 page (arcprize.org/competitions/2026) says ``all code and methods authored by the submitter must be made open source under a permissive public domain license (eg. CC0 or MIT-0)'', and that third-party code must be available under ``at least, an open source license which allows public sharing (eg. Apache-2.0, GPLv3)''. The same page says ``Participants must open source their solutions before receiving official private evaluation scores''; when this applies is not stated. Kaggle's rules for both competitions license a winning submission and its source code under CC-BY 4.0, a second licence text next to CC0 or MIT-0.
The Paper Track asks for a Writeup of at most 1,500 words with a cover image and an attached public notebook, on its Kaggle page. A PDF linked from the Writeup has no word limit \citep{kamradt2026forum748431}. A private Kaggle resource attached to a public Writeup ``will automatically be made public after the deadline'' (2026-11-09 23:59 UTC in the rules as fetched; arcprize.org says papers are due 2026-11-08). The Accuracy criterion matches the submission ID to the notebook. % N-E6s6.1 audit_cites-C8

\paragraph{The Qwen question.} Rule 5.a.1.a of the ARC-AGI-3 competition requires an open source system, model and weights as defined by the Open Source AI Definition checklist. Rule 5.a.3 says that input data or pretrained models with an incompatible licence need no open-source grant. The first sentence is absent from the Paper Track rules we fetched (a string search finds none there and one in the ARC-AGI-3 rules); the second appears in both. Forum threads 745079 and 745837 ask whether Qwen Community License weights are eligible. The hosts had not replied in either on 2026-10-13. We draw no conclusion and make no claim about eligibility. % N-E6s6.1

\paragraph{AI use.} The author directed the work and selected what went into submissions. Claude Code agents (Anthropic) wrote much of the code, the analyses and the text.\ownernote{Confirm or rewrite the two sentences before this note.}

\noindent The agent evaluated is the open-weight Qwen3.8-Flash-Next served locally; no external API is called at evaluation time. A scheduled job made the daily Kaggle submissions under the author's standing authorisation. Numbers in the paper come from recorded files and scripts in the repository (Appendix~\ref{app:regen}). AI systems are not authors. % H10
